\documentclass[10pt,a4paper]{amsart}
\usepackage[margin=19mm]{geometry}
\usepackage{amsmath,amssymb,mathtools}
\usepackage[scr=rsfs,cal=euler]{mathalfa}
\usepackage{xfrac}
\usepackage[unicode,hidelinks,bookmarks=false]{hyperref}
\newcommand{\ie}{i.e.\ }
\newcommand{\cf}{cf.\ }
\newcommand{\resp}{respectively\ }
\newcommand{\Subsection}{\subsection}
\providecommand{\nobreakdash}{\nobreak\hbox{-}}
\setlength{\emergencystretch}{2em}
\multlinegap0pt
\usepackage{amssymb}
\usepackage[shortlabels]{enumitem}
\newtheorem{lemma}{Lemma}
\title{Exceptional sets for compositions involving Euler's function, the sum-of-divisors function and Dedekind's function}
\author{Aimin Guo$^{*}$}
\date{Source: arXiv:2608.19972v2 (CC BY 4.0)}
\begin{document}
\maketitle

\noindent\emph{Source and licence.} Exceptional sets for compositions involving Euler's function, the sum-of-divisors function and Dedekind's function. Version arXiv:2608.19972v2 (CC BY 4.0), distributed by arXiv under the Creative Commons Attribution 4.0 International licence, http://creativecommons.org/licenses/by/4.0/, as recorded in the arXiv licence metadata for this exact version and shown on the abstract page https://arxiv.org/abs/2608.19972v2. Full licence notice: \texttt{licenses/arxiv-2608.19972-CC-BY-4.0.txt}.\par\smallskip

\noindent\textit{Editorial scope.} Counted unit: the article's lemma lem:second-moments (source0.tex lines 575--592). The article's complete original proof of it is delivered with it (source0.tex lines 594--895, one \texttt{proof} environment of 302 lines). No mathematics is written, repaired or re-derived: the statement and the proof are the article's own lines, in the article's own notation, and no retained line is substituted or re-flowed.  No further source-local context is needed: every cross-reference of every retained line resolves inside the two delivered spans.  Nothing was omitted from any retained span, so the package carries no omission record.  No retained line contains a citation, so the wrapper carries no bibliography and no bibliography entry.  No retained line contains a figure environment, an image-inclusion command or drawing code, so no image is delivered and the package declares no asset dependency.  Editorial formatting only, in addition: an AMS \texttt{amsart} preamble that restates the article's own declarations and macros used by the retained lines, namely the environments \texttt{lemma} on separate counters, together with the standard packages those lines need.  The retained environments print in this document as Lemma 1 in order of appearance, whereas the article prints the same items with the numbering of its own full document.  The complete native source of the article as distributed in the arXiv e-print for this exact version is bundled unchanged as \texttt{sources/208139/source0.tex}.
\begin{lemma}\label{lem:second-moments}
There exist absolute constants $C_{\psi},C_{\sigma}>0$ such that, for every
$x\geq 1$,
\[
\sum_{n\leq x}\left(\frac{\psi(n)}{n}\right)^2
\leq C_{\psi}x \qquad
and
\qquad
\sum_{n\leq x}\left(\frac{\sigma(n)}{n}\right)^2
\leq C_{\sigma}x.
\]
In particular,
\[
\sum_{n\leq x}\left(\frac{\psi(n)}{n}\right)^2\ll x,
\qquad
\sum_{n\leq x}\left(\frac{\sigma(n)}{n}\right)^2\ll x.
\]
\end{lemma}

\begin{proof}
We first consider the Dedekind function. Since
\[
\frac{\psi(n)}{n}
=
\prod_{p\mid n}\left(1+\frac1p\right),
\]
we have
\[
\left(\frac{\psi(n)}{n}\right)^2
=
\prod_{p\mid n}
\left(1+\frac2p+\frac1{p^2}\right).
\]
Define the multiplicative function \(a\) by
\[
a(1)=1,\qquad
a(p)=\frac2p+\frac1{p^2},
\qquad
a(p^\nu)=0\quad(\nu\geq2).
\]
Then
\[
\left(\frac{\psi(n)}{n}\right)^2
=
\sum_{d\mid n}a(d).
\]
Since \(a(d)\geq0\), it follows that
\[
\begin{aligned}
\sum_{n\leq x}\left(\frac{\psi(n)}{n}\right)^2
&=
\sum_{n\leq x}\sum_{d\mid n}a(d)\\
&=
\sum_{d\leq x}a(d)
\left\lfloor\frac{x}{d}\right\rfloor\\
&\leq
x\sum_{d\leq x}\frac{a(d)}{d}\\
&\leq
x\sum_{d=1}^{\infty}\frac{a(d)}{d}.
\end{aligned}
\]

We now show that the series on the right converges. For \(z\geq2\),
multiplicativity gives
\[
\sum_{\substack{d\geq1\\ p\mid d\Rightarrow p\leq z}}
\frac{a(d)}{d}
=
\prod_{p\leq z}
\left(1+\frac{a(p)}{p}\right)
=
\prod_{p\leq z}
\left(1+\frac2{p^2}+\frac1{p^3}\right).
\]
Put
\[
u_p:=\frac2{p^2}+\frac1{p^3}.
\]
Since \(p\geq2\), we have
\[
0<u_p
=
\frac2{p^2}+\frac1{p^3}
\leq
\frac3{p^2}.
\]
Moreover,
\[
\sum_p\frac1{p^2}
\leq
\sum_{n=2}^{\infty}\frac1{n^2}
=
\frac{\pi^2}{6}-1
<
\infty.
\]
Hence
\[
\sum_pu_p
\leq
3\sum_p\frac1{p^2}
<
\infty.
\]
Using the elementary inequality
\[
\log(1+t)\leq t
\qquad(t\geq0),
\]
we obtain
\[
0\leq
\sum_p\log(1+u_p)
\leq
\sum_pu_p
<
\infty.
\]
Therefore the Euler product
\[
\prod_p
\left(1+\frac2{p^2}+\frac1{p^3}\right)
\]
converges to a finite positive value.

Since \(a(d)\geq0\), letting \(z\to\infty\) in the finite Euler-product
identity above and using monotone convergence yields
\[
\sum_{d=1}^{\infty}\frac{a(d)}{d}
=
\prod_p
\left(1+\frac2{p^2}+\frac1{p^3}\right)
<\infty.
\]
We may therefore define
\[
C_\psi
:=
\sum_{d=1}^{\infty}\frac{a(d)}{d}
=
\prod_p
\left(1+\frac2{p^2}+\frac1{p^3}\right)
<\infty.
\]
Consequently,
\[
\sum_{n\leq x}
\left(\frac{\psi(n)}{n}\right)^2
\leq C_\psi x.
\]

We next consider the sum-of-divisors function. Put
\[
F(n):=\left(\frac{\sigma(n)}{n}\right)^2.
\]
Since \(\sigma(n)/n\) is multiplicative, so is \(F\). Define
\[
b:=F*\mu.
\]
Then \(b\) is multiplicative and
\[
F(n)=\sum_{d\mid n}b(d).
\]
For every prime power \(p^\nu\), \(\nu\geq1\), we have
\[
\frac{\sigma(p^\nu)}{p^\nu}
=
1+\frac1p+\cdots+\frac1{p^\nu}.
\]
Hence
\[
b(p^\nu)
=
\left(1+\frac1p+\cdots+\frac1{p^\nu}\right)^2
-
\left(1+\frac1p+\cdots+\frac1{p^{\nu-1}}\right)^2.
\]
In particular, \(b(p^\nu)\geq0\), and therefore \(b(d)\geq0\) for all \(d\).

Let
\[
S_\nu:=1+\frac1p+\cdots+\frac1{p^\nu}.
\]
Then
\[
b(p^\nu)
=
S_\nu^2-S_{\nu-1}^2
=
(S_\nu-S_{\nu-1})(S_\nu+S_{\nu-1}).
\]
Since
\[
S_\nu-S_{\nu-1}=\frac1{p^\nu}
\]
and
\[
S_\nu,S_{\nu-1}
\leq
\sum_{j=0}^{\infty}\frac1{p^j}
=
\frac{p}{p-1}
\leq2,
\]
we obtain
\[
0\leq b(p^\nu)\leq\frac4{p^\nu}.
\]

Since \(b(d)\geq0\), it follows that
\[
\begin{aligned}
\sum_{n\leq x}\left(\frac{\sigma(n)}{n}\right)^2
&=
\sum_{n\leq x}\sum_{d\mid n}b(d)\\
&=
\sum_{d\leq x}b(d)
\left\lfloor\frac{x}{d}\right\rfloor\\
&\leq
x\sum_{d\leq x}\frac{b(d)}{d}\\
&\leq
x\sum_{d=1}^{\infty}\frac{b(d)}{d}.
\end{aligned}
\]

We now show that the series on the right converges. For \(z\geq2\),
multiplicativity gives
\[
\sum_{\substack{d\geq1\\ p\mid d\Rightarrow p\leq z}}
\frac{b(d)}{d}
=
\prod_{p\leq z}
\left(
1+\sum_{\nu\geq1}\frac{b(p^\nu)}{p^\nu}
\right).
\]
Put
\[
v_p:=
\sum_{\nu\geq1}\frac{b(p^\nu)}{p^\nu}.
\]
By the estimate above,
\[
0\leq v_p
\leq
4\sum_{\nu\geq1}\frac1{p^{2\nu}}
=
\frac4{p^2-1}.
\]
Since \(p\geq2\),
\[
\frac1{p^2-1}\leq\frac2{p^2},
\]
and hence
\[
\sum_pv_p
\leq
8\sum_p\frac1{p^2}
\leq
8\sum_{n=2}^{\infty}\frac1{n^2}
=
8\left(\frac{\pi^2}{6}-1\right)
<
\infty.
\]
Again using
\[
\log(1+t)\leq t
\qquad(t\geq0),
\]
we obtain
\[
0\leq
\sum_p\log(1+v_p)
\leq
\sum_pv_p
<
\infty.
\]
Therefore the Euler product
\[
\prod_p(1+v_p)
=
\prod_p
\left(
1+\sum_{\nu\geq1}\frac{b(p^\nu)}{p^\nu}
\right)
\]
converges to a finite positive value.

Since \(b(d)\geq0\), letting \(z\to\infty\) in the finite Euler-product
identity above and using monotone convergence yields
\[
\sum_{d=1}^{\infty}\frac{b(d)}{d}
=
\prod_p
\left(
1+\sum_{\nu\geq1}\frac{b(p^\nu)}{p^\nu}
\right)
<\infty.
\]
We may therefore define
\[
C_\sigma
:=
\sum_{d=1}^{\infty}\frac{b(d)}{d}
=
\prod_p
\left(
1+\sum_{\nu\geq1}\frac{b(p^\nu)}{p^\nu}
\right)
<\infty.
\]
Consequently,
\[
\sum_{n\leq x}
\left(\frac{\sigma(n)}{n}\right)^2
\leq C_\sigma x.
\]
This completes the proof.
\end{proof}

\end{document}
